% !TeX root = ../教材题库_全部条目_转写工作稿.tex
% Source: XJ-B1-C04-HW-P149-Y0447
% 第4章 / 4.5 函数模型及其应用 / 教材 p.149 / PDF 第156页
% OCR未校对：公式与文字均可能有误；未完成数学TeX转写；禁止直接发给学生
% 原题及相邻原解截图（核对用）：../题目原图/XJ-B1-C04-HW-P149-Y0447-01.webp
% 本文件只在转写完成后作为题干使用；OCR 条目尚不是可靠题干。
\begin{bankitem}{习题4.5 第8题}{第4章 / 4.5 函数模型及其应用 / 教材 p.149 / PDF 第156页}{XJ-B1-C04-HW-P149-Y0447}{OCR未校对：公式与文字均可能有误；未完成数学TeX转写；禁止直接发给学生}
\begin{ocrdraft}
[教材p.149 / PDF第156页]
\par
8．几名中学生对“怎样烧开水最快最省煤气”的问题进行实验研究．当关着煤
\par
气的时候，煤气旋钮（下称旋钮）的位置定为0。，煤气开到最大时为90。．在0。、90。
\par
之间分成5等份，代表不同的煤气流量，它们分别是18。，36。，
\par
54。，72。，90。．实
\par
验数据如下表：
\par
旋钮位置
\par
18。
\par
36。
\par
54。
\par
72。
\par
烧开一壶水（3．75L)所需时间/min
\par
19
\par
16
\par
13
\par
12
\par
10
\par
烧开一壶水所需煤气量/m
\par
0．130
\par
0．122
\par
0．139
\par
0．149
\par
0．172
\par
试对表中数据做分析，并回答下列问题：
\par
（1)能否找到最省时间同时又最省气的烧水方法？
\par
（2）如果你急于用开水，最省时间的烧水方法是怎样的？
\par
（3）把旋钮位置作为自变量，烧开一壶水所需煤气量作为函数值，请你在平面
\par
直角坐标系内作出相应的图象．
\par
（4）从图象上看，是不是旋钮开得越小越省气？为什么？怎样才能比较省气？
\end{ocrdraft}
\end{bankitem}
